\section{Conditional Memory Evaluation}
\label{sec:method}

\subsection{The measured quantity}
Figure~\ref{fig:pipeline} shows the protocol end to end. We call the model that answers a future evaluation task the \emph{reader},
and reserve ``agent'' for the motivating setting. A \emph{memory unit} $u$ is one
page visit from a past episode; a memory $m$ is the Raw, Summary or Structured
text written from $u$. For a memory $m$ and a later, distinct future task $t$, the conditional utility of
$m$ for $t$ is the change in reward that adding $m$ causes, given the context
filler $F$, the position $p$ at which $m$ is inserted, and the seed $s$:
\begin{equation}
\label{eq:du}
  \Delta U(m, t \mid F, p, s) = R(t \mid F \cup_p m; s) - R(t \mid F; s).
\end{equation}
A positive value means the memory helped and a negative value that it hurt; a
zero difference indicates no measured change in reward, not identical answers
or absence of useful information. We call these labels \emph{helped},
\emph{unchanged} and \emph{hurt}; each holds under its evaluation condition. Both arms fix the reader model, serving
cluster, decoding, prompt scaffold, output-token budget, and scoring function
and version; they differ only in whether $m$ is present. We call them the
\emph{with-target-memory} and \emph{without-target-memory} arms (\emph{with $m$} and
\emph{without $m$}); both contain the filler. Reward stays continuous
wherever the benchmark gives partial credit; in the screened run it is
binary. The per-cell estimate is the mean over replicates. We write $\Delta
U(m,t)$ when $F$, $p$ and $s$ are fixed, and the conditioning matters: a single
inserted item can change the answer for reasons unrelated to its content. The unit of analysis is the
\emph{cell}, a memory--task pair $(m,t)$ under one evaluation condition; a memory
has no single utility here, and its result is its profile across the tasks it
faced. We compare at three levels: the with-$m$ versus without-$m$ reward
difference of each pair; within each memory, the mean difference on
answer-bearing task pairs minus that on matched-control task pairs (the
headline contrast); and representations written from the same memory unit.

\paragraph{Identical filler, frozen memories.}
Both arms place the same controlled filler in the reader's context, and the
\textsc{with} arm adds $m$. The filler is $\lfFillerK$ memory units of
$\lfFillerLo$--$\lfFillerHi$ tokens each, drawn from $\lfFillerTraj$
trajectories disjoint from every trajectory that hosts a bearing memory, so it
never contains another unit that bears on the task. The \textsc{without} arm
is therefore not an empty-context arm. On LongMemEval's first version~\cite{lmev1} the oracle context is entirely
evidence, so ablating ``the memory'' there would delete the whole context. The memory set is written once, frozen and content-hashed before any
pairing exists, and no arm writes, updates or deletes a memory. The arms
differ by exactly one item from that frozen set, so the divergent-state
hazard raised by Memory-R2~\cite{memoryrtwo} does not arise. The position of
the added memory is not neutral: moving only that position changed the answer
in 8 of 12 cells of a pilot, and $\soPosChgPct\%$ of labels at scale
(Section~\ref{sec:limitations}). In the main runs we hold position fixed across replicates, vary it
across cells, and record it on every row (\supp{sec:position}). Every
contrast we report is therefore conditional on the recorded positions;
position variance is not inside our intervals.

\paragraph{The writer never sees the future task.}
Memory formation is a function of the memory unit alone, $m = W(u)$ and
never $m = W(u,t)$. Without this invariant, $\Delta U > 0$ would only show
that a model can be told an answer and copy it forward. We enforce it
structurally rather than by convention: the writer takes a segment and
nothing else, emits a closed field set, and has a single interpolation point.
Each of these barriers carries a negative control that plants the defect and requires
the check to fire (\supp{sec:taskblind}).

\subsection{The answer-bearing screen is an instrument, not a retriever}
\label{sec:screen}
To decide which (memory, task) pairs are worth executing, we use an
\emph{answer-bearing screen}. For each pair it reads three inputs: the memory's text,
the task's \textbf{gold answer}, and the task's evaluation specification. It
never reads the question, a reward, or any model output, so the strata it
defines are not selected on the outcome. A pair is \emph{answer-bearing} (\textsc{bearing}) when the
benchmark's own scorer would mark the memory text correct against the gold
answer \emph{and} the test singles this memory out of its domain pool, firing
on at most eight of the 865--1{,}024 units available to that task.

Answer-bearing is a screening label based on answer containment, not a
measured positive-utility label: paired evaluation measures the effect, and
screen-selected pairs can still fail while pairs outside the screen can still
help. Our sparsity result is a statement about how answers are distributed
over memories. Because it consumes the gold answer, the screen
cannot be deployed, and we do not propose it as a retriever. Its role is to define, with the gold answer, the
answer-bearing pairs; these screen-defined labels, not exhaustive ground truth
for useful memories, are what we evaluate retrieval against (Section~\ref{sec:similarity}). Every rate we
report under the screen is conditional on a screened design.

A screen selected on the gold answer invites a second objection: perhaps the
effect is an artifact of selecting pairs with the answer in hand. We tested
this with a gold-free screen, pre-registered before execution: the same model
as the reader (\lfGfJudge), asked whether the memory contains the specific
information the question asks about, without the gold answer. On $\gfCells$
cells the gold screen \emph{never} selected, the gold-free picks move the
reader from $\gfWithoutPct\%$ correct without the memory to $\gfWithPct\%$ with
it, while $\lfGfCtlCells$ matched-control pairs move from
$\gfCtlWithoutPct\%$ to $\gfCtlWithPct\%$. The contrast is $\rd{\gfDiff}$, with a
task-clustered 95\% interval of $[\lfGfTaskLo, \lfGfTaskHi]$ over
$\lfGfTaskN$ tasks and an episode-clustered one of $[\lfGfEpLo, \lfGfEpHi]$
over $\lfGfEpN$ episodes. This answers the selection objection, not the
containment one; because the judge is the reader's own model, its picks may
favor pairs that model finds easy. Its picks overlap the gold screen little
(\supp{sec:goldfree}) yet carry an effect of similar size, so benefits also
occur outside the gold screen. A second, lexical gold-free screen did not
locate answer-bearing pairs (\supp{sec:goldfree}).

\subsection{Three representations of one experience}
All three representations derive from the \emph{same} memory unit, so a
downstream difference is attributable to how the segment is written rather
than to which content was selected. \textbf{Raw} is deterministic and uses no
model: the memory is the segment itself. \textbf{Summary} is one writer call
per segment at pinned decoding and prompt version. \textbf{Structured} is one
writer call that emits a validated schema (memory type, trigger, action,
outcome, lesson, scope). A malformed structure is recorded as a failure and
never repaired.

Formation cost belongs to the representation. Raw makes zero writer calls;
Summary and Structured make one each per segment. Omitting that cost flatters
the model-written representations, so we charge each its own, including
writes that failed after consuming a call.

\subsection{Noise, labels and cost}
\label{sec:metrics}
Noise is measured, not assumed. Replicates are separate requests
($\lfScrCacheHits$ of $\lfScrRealCalls$ calls served from cache), and within a
session they agree on all 1{,}157 primary cells, so the execution floor is zero
and a cell's label (helped, unchanged or hurt) is the sign of $\Delta U$; the pre-registered floors and
label rules are in \supp{sec:floors}. Truncated completions are re-issued, not
scored. Intervals are over memories, tasks or episodes, never replicates.

A memory shows a \textbf{sign flip} when it helps one task and hurts another. A flip rate is descriptive, not evidence of
memory\,$\times$\,task interaction: under a purely additive model
$\Delta U = \mu + a_m + b_t$, simulated over 20 trials of 200 memories
$\times$ 8 tasks, flip rates range from 53.5\% to 92.0\% with no interaction
present (\supp{sec:floors}). A flip rate also grows with the number of cells
per memory, so it can only be compared across groups at matched cell counts.

\paragraph{The perturbation floor.}
The noise that matters most here is not seed variation but the reader's
response to inserting \emph{any} item: position alone changes
$\soPosChgPct\%$ of per-cell labels (Section~\ref{sec:limitations}). We ran no placebo arm of length-matched irrelevant content,
so per-cell labels include such perturbations. The random-task stratum bounds their size: a memory paired with a randomly
selected task helps in $\raHelpRateRand\%$ of the pairs where the reader was
wrong without it and hurts in $\raHurtRateRand\%$ of those where it was right.
These rates are an upper bound on pure perturbation: the memories are real
same-domain pages, and where they contain the answer string they help more
often ($\rvRandFireHelp\%$ against $\rvRandSilentHelp\%$). Irrelevant
context can distract a reader~\cite{distracted} or even help it~\cite{powernoise}. We therefore report effects as paired contrasts between strata, which reduce
sensitivity to content-independent perturbation when its mean effect is
comparable across the contrasted strata, and never as counts of labeled cells;
without a placebo arm we do not claim to isolate every such effect.

\paragraph{Pre-registered and exploratory analyses.}
The screened design, both gold-free screens, the length control, the
second-reader study, the writer\,$\times$\,reader factorial, the reversal rate
and the variance decomposition were pre-registered; \supp{sec:prereg} lists
their anchors and every deviation. All other analyses, including the three
follow-up runs, are exploratory, and their intervals are not adjusted for
multiplicity.

Every measurement is paired with its cost in seven families, from storage to
formation; utility per exposed token is reported beside utility per token
charged, including the writer's (\supp{sec:reuse}).

